% q0050.tex — Who ate my cake?
\question{
  id        = {0050},
  title     = {Who ate my cake?},
  source    = {OBECON},
  year      = {2026},
  round     = {Seletiva IEO -- Phase C},
  language  = {English},
  topic     = {Macro},
  subtopic  = {Inflation / Shrinkflation},
  type      = {objective},
  statement = {%
    Argentiland has 50\% monthly inflation; Brazilkistan has 3\% annual
    inflation. In Argentiland, companies keep product sizes stable but raise
    prices regularly. In Brazilkistan, companies reduce product sizes while
    keeping nominal prices stable (shrinkflation). What best explains this
    difference in firm behaviour?
  },
  choices   = {%
    \item Brazilkistan consumers are less likely to accept price increases.
    \item Argentiland consumers anticipate frequent price increases, so
          companies can raise prices without changing product sizes.
    \item Argentiland companies face higher production costs.
    \item Higher inflation makes Argentiland companies prefer price increases
          to shrinkflation.
  },
  answer    = {B},
  solution  = {%
    In high-inflation environments like Argentiland, price increases are
    anticipated and normalised, greatly reducing consumer resistance to nominal
    price changes. In moderate-inflation economies like Brazilkistan, consumers
    are far more sensitive to visible price increases, making shrinkflation --- a
    covert real price increase --- the preferred strategy for maintaining revenues
    without triggering consumer backlash. (B) is correct.
  },
}
